Memory-augmented LLM agents must track not only \emph{what} happened but \emph{when} it happened and \emph{when they learned about it}, a dual-timestamp (bitemporal) distinction critical for real-world applications such as medical records, financial systems, and news analysis.
Existing benchmarks either ignore timestamps entirely or treat time as a single dimension, leaving a gap in evaluating this capability.

We introduce \textbf{BiTempQA}, a diagnostic benchmark of 415 Chinese question-answer pairs (84 held-out test, all sourced from real-world data) across 10 scenario types and 3 question types at two difficulty levels, where each memory entry carries explicit \texttt{event\_time} (when something occurred) and \texttt{record\_time} (when the system stored it) annotations.
Over half (56.0\%) of test questions require reasoning about both timestamps simultaneously.
The dataset combines ChronoQA news data (200 scenarios) with 15 manually curated real-world temporal event scenarios and 118 LLM-generated scenarios covering rare types (64.6\% real-source overall).
We evaluate seven memory backends (FAISS, BM25, Naive RAG, Simple KG, ChromaDB, Mem0, Graphiti) through a unified Retrieve$\to$LLM Generate$\to$LLM Judge pipeline with controlled evaluation parameters (same embedding model, same top-$k$, same generator).
We find that: (1) the top four lightweight systems are statistically indistinguishable (81.0--82.1\%, $p{>}0.3$), suggesting that retrieval design matters less than answer generation quality for short-context bitemporal scenarios; (2) ChromaDB's metadata filtering approach fails catastrophically (22.6\%) due to heterogeneous temporal formats in real-world data; (3) production memory systems (Mem0, Graphiti) dramatically underperform lightweight baselines, revealing that aggressive memory compression destroys bitemporal detail; and (4) real-world and LLM-generated scenarios yield comparable accuracy, validating our data augmentation approach.
BiTempQA is released as an open diagnostic tool for developing temporally-aware memory systems.
